% ROZDZIAŁ 7 — do włączenia w main.tex
\chapter{Konstrukcja całki Lebesgue'a}
\label{ch:calkowanie-lebesgue}

\section{Od mierzenia zbiorów do mierzenia wartości funkcji}

W klasycznej sumie Riemanna dzielimy \emph{dziedzinę} na małe
przedziały. Konstrukcja Lebesgue'a rozpoczyna się od podziału
\emph{zbioru wartości}: dla funkcji prostej $s$ mierzymy zbiory,
na których $s$ przyjmuje kolejne wartości, i sumujemy odpowiednie
iloczyny. Pozwala to całkować na dowolnych przestrzeniach z miarą,
bez porządku na dziedzinie i bez pojęcia przedziału.

W całym rozdziale $(X,\mathcal F,\mu)$ jest dowolną przestrzenią
z miarą. Nie zakładamy, że $\mu(X)<\infty$, ani nawet
$\sigma$-skończoności. Jeżeli nie zaznaczamy inaczej,
funkcje są mierzalne względem $\mathcal F$.
Oznaczenia $\overline{\mathbb R}=[-\infty,\infty]$,
$f^+=\max(f,0)$ i $f^-=\max(-f,0)$ mają zwykły sens;
$f=f^+-f^-$ w punktach, gdzie $f$ jest skończona.
Dla nieujemnych wielkości z $[0,\infty]$ umawiamy się,
że $0\cdot\infty=0$, również przy wyznaczaniu całki
nieujemnej funkcji prostej. Nie dopuszczamy odejmowania
$\infty-\infty$.

\section{Całka nieujemnej funkcji prostej}

\begin{definicja}[Całka funkcji prostej]
\label{def:integral-simple}
Niech $s:X\to[0,\infty)$ będzie mierzalną funkcją prostą.
Jeżeli $a_1,\ldots,a_m$ są wszystkimi różnymi wartościami $s$,
a $A_j=\{x\in X:s(x)=a_j\}$, określamy
\begin{equation}\label{eq:simple-integral}
\int_X s(x)\,\mathrm d\mu(x)
:=\sum_{j=1}^{m}a_j\mu(A_j)\in[0,\infty].
\end{equation}
Zbiory $A_j$ są mierzalne i tworzą rozłączny podział $X$.
Przy $a_j=0$ składnik jest równy zeru, niezależnie od
$\mu(A_j)$. Wartość całki może być nieskończona.
\end{definicja}

\begin{lemat}[Niezależność od reprezentacji]
\label{lem:simple-refinement}
Jeżeli $s=\sum_{i=1}^p a_i\chi_{A_i}$, gdzie
$a_i\ge0$, a zbiory $A_i\in\mathcal F$ tworzą
rozłączny podział $X$ (dopuszczamy powtarzające się
współczynniki), to
$\int_Xs\,\mathrm d\mu(x)=\sum_{i=1}^pa_i\mu(A_i)$.
Wartość ta nie zależy od wybranego podziału ani od
powtarzania współczynników.
\end{lemat}
\begin{proof}
Niech $a_1',\ldots,a_r'$ będą różnymi wartościami
funkcji $s$, a $C_j=\{s=a_j'\}$. Dla każdego $j$
zbiór $C_j$ jest rozłączną sumą tych $A_i$, dla których
$a_i=a_j'$. Zatem, na mocy skończonej addytywności miary,
$a_j'\mu(C_j)=\sum_{i:a_i=a_j'}a_i\mu(A_i)$.
Gdy $a_j'=0$, obie strony są zerem według przyjętej
konwencji; dla $a_j'>0$ równość wynika z rozdzielności
mnożenia w $[0,\infty]$. Sumując po $j$, otrzymujemy
tezę. Dla dwóch podziałów można również zastosować
wspólne rozdrobnienie $A_i\cap B_j$; dowodzi to
niezależności bez zakładania skończoności miary.
\end{proof}

\begin{twierdzenie}[Rachunek dla funkcji prostych]
\label{thm:simple-properties}
Niech $s,t:X\to[0,\infty)$ będą mierzalnymi funkcjami
prostymi, a $c\in[0,\infty)$. Wtedy:
\begin{enumerate}[label=\textup{(\roman*)}]
\item $\int_X0\,\mathrm d\mu(x)=0$ oraz
$\int_X\chi_A\,\mathrm d\mu(x)=\mu(A)$ dla każdego
$A\in\mathcal F$;
\item jeżeli $s\le t$, to
$\int_Xs\,\mathrm d\mu(x)\le\int_Xt\,\mathrm d\mu(x)$;
\item $\int_X(cs)\,\mathrm d\mu(x)=c\int_Xs\,\mathrm d\mu(x)$;
\item $\int_X(s+t)\,\mathrm d\mu(x)=
\int_Xs\,\mathrm d\mu(x)+\int_Xt\,\mathrm d\mu(x)$.
\end{enumerate}
W punktach (iii) i (iv) wszystkie działania odbywają się
w $[0,\infty]$ z konwencją $0\cdot\infty=0$.
\end{twierdzenie}
\begin{proof}
Pierwsza część wynika wprost z definicji. Niech
$s=\sum_i a_i\chi_{A_i}$ i
$t=\sum_j b_j\chi_{B_j}$ będą reprezentacjami
na rozłącznych podziałach $X$. Zbiory
$D_{ij}=A_i\cap B_j$ tworzą wspólne rozdrobnienie.
Na każdym $D_{ij}$ obie funkcje są stałe, dlatego
\[
\int_Xs\,\mathrm d\mu(x)=\sum_{i,j}a_i\mu(D_{ij}),\qquad
\int_Xt\,\mathrm d\mu(x)=\sum_{i,j}b_j\mu(D_{ij}).
\]
Jeśli $s\le t$, to dla każdego niepustego $D_{ij}$
mamy $a_i\le b_j$, więc każdy składnik pierwszej sumy
jest nie większy od odpowiedniego składnika drugiej.
Skalowanie i addytywność wynikają z równości
$(ca_i)\mu(D_{ij})=c(a_i\mu(D_{ij}))$ oraz
$(a_i+b_j)\mu(D_{ij})=a_i\mu(D_{ij})+b_j\mu(D_{ij})$.
Żadna operacja nie wymaga odejmowania nieskończoności.
\end{proof}

\begin{przyklad}[Mierzalna funkcja prosta na przestrzeni nieskończonej]
Na $(\mathbb R,\mathcal L_1,\lambda_1)$ funkcja
$s=3\chi_{[0,2]}+5\chi_{[4,5]}$ ma całkę
$3\cdot2+5\cdot1=11$. Natomiast
$t=\chi_{[0,\infty)}$ ma całkę $+\infty$.
Obie funkcje są proste i mierzalne, chociaż druga
nie jest całkowalna w sensie skończonej całki.
\end{przyklad}

\section{Definicja całki dowolnej funkcji nieujemnej}

\begin{definicja}[Całka Lebesgue'a funkcji nieujemnej]
\label{def:integral-positive}
Niech $f:X\to[0,\infty]$ będzie mierzalna. Definiujemy
\begin{equation}\label{eq:positive-integral}
\int_X f(x)\,\mathrm d\mu(x)
:=\sup\left\{\int_X s(x)\,\mathrm d\mu(x):
\begin{array}{l}
 s:X\to[0,\infty)\text{ jest mierzalna i prosta},\\[-2pt]
 0\le s(x)\le f(x)\text{ dla wszystkich }x\in X
\end{array}\right\}.
\end{equation}
Supremum należy do $[0,\infty]$. Rodzina funkcji
prostych w nawiasie jest niepusta, ponieważ zawiera zero.
\end{definicja}

\begin{uwaga}[Co dokładnie znaczy ta konstrukcja?]
Mierzalność funkcji $f$ nie wymusza skończoności całki.
Nie wybieramy też żadnego konkretnego przybliżenia:
rozpatrujemy \emph{wszystkie} nieujemne funkcje proste
nieprzekraczające $f$. Twierdzenie
\ref{thm:simple-up} zapewnia, że istnieją również
jawne, punktowo niemalejące przybliżenia $s_n\uparrow f$.
Dopiero twierdzenie Beppo Leviego w rozdziale
\ref{ch:twierdzenia-graniczne} wykaże, że całkę można
obliczać jako granicę całek takiej sekwencji.
\end{uwaga}

\begin{twierdzenie}[Pierwsze własności całki nieujemnej]
\label{thm:nonnegative-first-properties}
Niech $f,g:X\to[0,\infty]$ będą mierzalne,
$c\in[0,\infty)$ i $A\in\mathcal F$. Wtedy:
\begin{enumerate}[label=\textup{(\roman*)}]
\item gdy $f\le g$, zachodzi
$\int_X f\,\mathrm d\mu(x)\le\int_Xg\,\mathrm d\mu(x)$;
\item $\int_X cf\,\mathrm d\mu(x)=c\int_Xf\,\mathrm d\mu(x)$;
\item dla funkcji prostej $f$ definicje
\ref{def:integral-simple} i \ref{def:integral-positive}
są zgodne;
\item $\int_X f\chi_A\,\mathrm d\mu(x)$ zależy
wyłącznie od wartości $f$ na $A$ i dla
$A\subseteq B$, $A,B\in\mathcal F$, zachodzi
$\int_X f\chi_A\,\mathrm d\mu(x)\le
\int_Xf\chi_B\,\mathrm d\mu(x)$;
\item jeśli $\mu(A)=0$, to
$\int_X f\chi_A\,\mathrm d\mu(x)=0$,
nawet gdy $f(x)=+\infty$ w punktach $A$.
\end{enumerate}
\end{twierdzenie}
\begin{proof}
W (i) każda funkcja prosta $s\le f$ spełnia
również $s\le g$, więc supremum dla $f$
nie przekracza supremum dla $g$.
Dla $c>0$ podstawienie $t=cs$ daje bijekcję
między nieujemnymi prostymi minorantami $f$ i $cf$;
na mocy twierdzenia \ref{thm:simple-properties}
całki skalują się przez $c$. Przy $c=0$ obie
strony wynoszą zero według umowy $0\cdot\infty=0$.

Jeżeli $f$ jest prosta, sama należy do rodziny
minorant w \eqref{eq:positive-integral}. Dla każdego
$s\le f$ monotoniczność całki prostej daje
$\int_Xs\,\mathrm d\mu(x)\le\int_Xf\,\mathrm d\mu(x)$;
stąd zgodność obu definicji.
Część (iv) wynika z punktowej równości funkcji
$f\chi_A$ przy ustalonych wartościach na $A$
oraz z (i). Aby wykazać (v), ustalmy dowolną
nieujemną funkcję prostą $s\le f\chi_A$.
Poza $A$ mamy $s=0$; skoro $\mu(A)=0$,
każdy jej dodatni poziom ma miarę zero, więc
$\int_Xs\,\mathrm d\mu(x)=0$.
Biorąc supremum, otrzymujemy zero.
\end{proof}

\begin{definicja}[Całka po podzbiorze]
\label{def:integral-over-set}
Dla $A\in\mathcal F$ i mierzalnej $f:X\to[0,\infty]$
określamy
$\int_A f(x)\,\mathrm d\mu(x):=
\int_X f(x)\chi_A(x)\,\mathrm d\mu(x)$,
gdzie iloczyn $f\chi_A$ rozumiemy jako funkcję
równą $f$ na $A$ i zero na $X\setminus A$.
To ostatnie określenie usuwa dwuznaczność zapisu
$0\cdot\infty$. Równoważnie całkujemy funkcję
$f|_A$ względem ograniczenia miary $\mu|_A$
na $(A,\mathcal F|_A)$.
\end{definicja}

\begin{lemat}[Zgodność z ograniczeniem miary]
\label{lem:restriction-integral}
Dla dowolnych $A\in\mathcal F$ i mierzalnej
$f:X\to[0,\infty]$ zachodzi
\[
\int_A f(x)\,\mathrm d\mu(x)
=\int_A f|_A(x)\,\mathrm d(\mu|_A)(x).
\]
\end{lemat}
\begin{proof}
Dla funkcji prostej zachodzi to z definicji i
$\mu|_A(A\cap B)=\mu(A\cap B)$.
Jeżeli $t:A\to[0,\infty)$ jest prostą mierzalną
minorantą $f|_A$, to przedłużenie jej zerem poza $A$
jest prostą mierzalną funkcją na $X$, mniejszą od
$f\chi_A$. Odwrotnie, ograniczenie każdej
prostej minoranty $f\chi_A$ do $A$ jest
prostą minorantą $f|_A$. Równość całek prostych
i supremów daje tezę.
\end{proof}

\begin{twierdzenie}[Nierówność Markowa dla całki]
\label{thm:markov-nonnegative}
Niech $f:X\to[0,\infty]$ będzie mierzalna i $a>0$.
Wówczas
\begin{equation}\label{eq:markov}
\mu(\{x\in X:f(x)\ge a\})
\le \frac1a\int_X f(x)\,\mathrm d\mu(x).
\end{equation}
Jeżeli $\int_Xf\,\mathrm d\mu(x)<\infty$, to $f$ jest
skończona $\mu$-prawie wszędzie. Ponadto
$\int_Xf\,\mathrm d\mu(x)=0$ wtedy i tylko wtedy,
gdy $f=0$ $\mu$-prawie wszędzie.
\end{twierdzenie}
\begin{proof}
Dla $E_a=\{f\ge a\}$ prosta funkcja
$a\chi_{E_a}$ nie przekracza $f$. Zatem
$a\mu(E_a)=\int_X a\chi_{E_a}\,\mathrm d\mu(x)
\le\int_Xf\,\mathrm d\mu(x)$, co daje
\eqref{eq:markov}. Jeśli ostatnia całka jest skończona,
to $\mu(\{f=\infty\})\le\mu(\{f\ge n\})
\le n^{-1}\int_Xf\,\mathrm d\mu(x)\to0$.
Jeżeli całka wynosi zero, to
$\mu(\{f\ge 1/n\})=0$ dla wszystkich $n$, a zbiór
$\{f>0\}=\bigcup_{n\ge1}\{f\ge1/n\}$ jest zerowy.
Odwrotnie, jeśli $f=0$ poza mierzalnym zbiorem zerowym,
to część (v) twierdzenia
\ref{thm:nonnegative-first-properties} daje $\int_X f=0$.
\end{proof}

\begin{przyklad}[Zero prawie wszędzie, mimo dużych wartości]
Na $\mathbb R$ z miarą Lebesgue'a funkcja
$f=+\infty\cdot\chi_{\mathbb Q}$, rozumiana jako
$f(x)=+\infty$ dla $x\in\mathbb Q$ i $0$ dla
$x\notin\mathbb Q$, jest mierzalna i ma całkę zero.
To pokazuje, dlaczego nie wystarczy traktować
całki jako zwykłego sumowania wartości punktowych.
\end{przyklad}

\section{Funkcje ze znakiem i całkowalność}

\begin{twierdzenie}[Niezmienniczość całki przy zmianie na zbiorze zerowym]
\label{thm:ae-integral-invariance}
Niech $f,g:X\to[0,\infty]$ będą mierzalne i
$f=g$ $\mu$-prawie wszędzie. Wówczas
$\int_Xf\,\mathrm d\mu(x)=\int_Xg\,\mathrm d\mu(x)$.
Jeśli $f,g:X\to\overline{\mathbb R}$ mają określone
całki ze znakiem i są równe prawie wszędzie, ich
całki również są równe.
\end{twierdzenie}
\begin{proof}
Wybierzmy $N\in\mathcal F$ o $\mu(N)=0$ takie,
że $f=g$ poza $N$. Dla dowolnej prostej minoranty
$0\le s\le f$ funkcja $s\chi_{X\setminus N}$
jest prostą minorantą $g$, a
\[
\int_Xs\,\mathrm d\mu(x)
=\int_Xs\chi_{X\setminus N}\,\mathrm d\mu(x),
\]
ponieważ składniki odpowiadające częściom poziomic
zawartym w $N$ mają miarę zero. Stąd
$\int_X f\le\int_X g$ po przejściu do supremum.
Zamiana ról $f,g$ daje nierówność odwrotną.
W przypadku funkcji ze znakiem stosujemy część
pierwszą osobno do ich części dodatnich i ujemnych.
Równość wartości wyklucza możliwość, by tylko jedna
z odpowiadających całek była nieokreślona.
\end{proof}

\begin{definicja}[Części dodatnia i ujemna]
\label{def:positive-negative-parts}
Dla mierzalnej $f:X\to\overline{\mathbb R}$ definiujemy
$f^+(x)=\max\{f(x),0\}$ i
$f^-(x)=\max\{-f(x),0\}$. Są to mierzalne funkcje
$[0,\infty]$-wartościowe, a $|f|=f^++f^-$.
Jeśli co najmniej jedna z liczb
$\int_Xf^+\,\mathrm d\mu(x)$,
$\int_Xf^-\,\mathrm d\mu(x)$ jest skończona, definiujemy
\begin{equation}\label{eq:signed-integral}
\int_X f(x)\,\mathrm d\mu(x)
:=\int_Xf^+(x)\,\mathrm d\mu(x)
-\int_Xf^-(x)\,\mathrm d\mu(x).
\end{equation}
Jeśli obie liczby są nieskończone, całka funkcji $f$
\emph{nie jest określona}. Funkcję $f$ nazywamy
\emph{(absolutnie) całkowalną}, jeżeli
$\int_X|f|\,\mathrm d\mu(x)<\infty$.
\end{definicja}

\begin{uwaga}[Dlaczego nie używamy odejmowania nieskończoności?]
Niech $X=\mathbb R$ z miarą Lebesgue'a i
$f(x)=x$. Wtedy zarówno całka $f^+$, jak i $f^-$
wynosi $\infty$, więc całka $f$ na całej $\mathbb R$
nie jest określona. Symetria i formalne równanie
"$\infty-\infty=0$" nie zastępują definicji.
Wartość główna Cauchy'ego, jeśli jest rozważana,
stanowi odrębne pojęcie wymagające osobnej definicji.
\end{uwaga}

\begin{twierdzenie}[Kryterium absolutnej całkowalności]
\label{thm:integrability-criterion}
Dla mierzalnej $f:X\to\overline{\mathbb R}$ następujące
warunki są równoważne:
\begin{enumerate}[label=\textup{(\roman*)}]
\item $\int_X|f|\,\mathrm d\mu(x)<\infty$;
\item $\int_Xf^+\,\mathrm d\mu(x)<\infty$ oraz
$\int_Xf^-\,\mathrm d\mu(x)<\infty$.
\end{enumerate}
W takim przypadku $f$ jest skończona prawie wszędzie,
całka $\int_Xf\,\mathrm d\mu(x)$ jest rzeczywista i
$|\int_Xf\,\mathrm d\mu(x)|\le\int_X|f|\,\mathrm d\mu(x)$.
\end{twierdzenie}
\begin{proof}
Ponieważ $0\le f^\pm\le|f|$, monotoniczność całki
daje (i)$\Rightarrow$(ii). Dla implikacji odwrotnej
nie możemy jeszcze korzystać bez dowodu z addytywności
całki dla dowolnych funkcji nieujemnych. Do pełnej addytywności wrócimy w rozdziale
\ref{ch:twierdzenia-graniczne}. Tu wykorzystamy
wyłącznie rozdzielność nośników części dodatniej i ujemnej.
Bezpośredni argument jest następujący: jeśli
$s\le f^++f^-$ jest nieujemna i prosta, to
dzielimy dziedzinę na mierzalne zbiory
$P=\{f\ge0\}$ i $M=\{f<0\}$. Funkcje
$s\chi_P\le f^+$ i $s\chi_M\le f^-$
są proste, a
$\int_Xs\,\mathrm d\mu(x)=
\int_Xs\chi_P\,\mathrm d\mu(x)+
\int_Xs\chi_M\,\mathrm d\mu(x)$
na mocy twierdzenia \ref{thm:simple-properties}.
Przechodząc do supremum po $s\le|f|=f^++f^-$,
otrzymujemy
$\int_X|f|\,\mathrm d\mu(x)\le
\int_Xf^+\,\mathrm d\mu(x)+
\int_Xf^-\,\mathrm d\mu(x)<\infty$.
Z twierdzenia \ref{thm:markov-nonnegative} otrzymujemy
skończoność $|f|$ prawie wszędzie. Wreszcie, dla dowolnych prostych minorant
$0\le u\le f^+$ i $0\le v\le f^-$ zachodzi
$u+v\le|f|$. Stąd
$\int_Xu\,\mathrm d\mu(x)+\int_Xv\,\mathrm d\mu(x)
=\int_X(u+v)\,\mathrm d\mu(x)
\le\int_X|f|\,\mathrm d\mu(x)$.
Przechodząc niezależnie do supremum po $u$ i $v$,
uzyskujemy także nierówność
$\int_Xf^+\,\mathrm d\mu(x)+\int_Xf^-\,\mathrm d\mu(x)
\le\int_X|f|\,\mathrm d\mu(x)$.
Z obydwu nierówności wynika równość tych liczb.
W szczególności $|\int_Xf\,\mathrm d\mu(x)|
\le\int_Xf^+\,\mathrm d\mu(x)+
\int_Xf^-\,\mathrm d\mu(x)=\int_X|f|\,\mathrm d\mu(x)$.
\end{proof}

\begin{definicja}[Całka funkcji całkowalnej po zbiorze]
\label{def:signed-integral-over-set-7}
Niech $A\in\mathcal F$, a $f:X\to\mathbb R$ będzie
mierzalna oraz całkowalna, tj.
$\int_X|f(x)|\,\mathrm d\mu(x)<\infty$.
Definiujemy
\[
\int_A f(x)\,\mathrm d\mu(x)
:=\int_X\widetilde f_A(x)\,\mathrm d\mu(x),
\qquad
\widetilde f_A(x)=
\begin{cases}f(x),&x\in A,\\0,&x\notin A.\end{cases}
\]
Funkcja $\widetilde f_A$ jest mierzalna oraz
$|\widetilde f_A|\le|f|$, więc całka po prawej
jest skończona. Zgodnie z definicją całki ze znakiem
również
$\int_A f=\int_A f^+-\int_A f^-$;
obie całki są skończone. Definicja nie wymaga
żadnej konwencji dla iloczynu $0\cdot\infty$.
\end{definicja}

\begin{uwaga}[Szczegół przy wartości $+\infty$]
W dowodzie wprowadziliśmy rozbicie
$P=\{f\ge0\}$ i $M=\{f<0\}$.
Dla $f(x)=+\infty$ obowiązuje $x\in P$,
a dla $f(x)=-\infty$ obowiązuje $x\in M$.
Nie wykonujemy w tych punktach odejmowania
$\infty-\infty$: obie części dodatnia i ujemna
nie są jednocześnie dodatnie w żadnym punkcie.
\end{uwaga}

\begin{przyklad}[Miary dyskretne i sumy szeregów]
Niech $X=\mathbb N$, $\mathcal F=\mathcal P(\mathbb N)$
i $\mu(A)=\sum_{n\in A}w_n$, gdzie $w_n\ge0$.
Dla funkcji $f:\mathbb N\to[0,\infty]$ całka
jest równa zwykłej sumie nieujemnego szeregu
$\int_{\mathbb N} f(n)\,\mathrm d\mu(n)=
\sum_{n=1}^{\infty}f(n)w_n$.
Istotnie, suma po pierwszych $N$ punktach jest
całką prostej funkcji nieprzekraczającej $f$,
więc prawa strona nie przekracza całki.
Z drugiej strony, dla dowolnej prostej $s\le f$
przeliczalna addytywność miary pozwala zapisać
$\int_X s\,\mathrm d\mu(x)=
\sum_{n=1}^{\infty}s(n)w_n\le\sum_{n=1}^{\infty}f(n)w_n$.
Dla rzeczywistej $f$ warunek absolutnej całkowalności
jest równoważny
$\sum_n|f(n)|w_n<\infty$.
W szczególności miara licząca odpowiada $w_n=1$.
\end{przyklad}

\begin{przyklad}[Całka może istnieć jako $+\infty$]
Na $(0,1]$ z miarą Lebesgue'a połóżmy $f(x)=1/x$.
Dla $j=0,\ldots,n-1$ na przedziale
$I_j=(2^{-(j+1)},2^{-j}]$ mamy $f(x)\ge2^j$.
Funkcja prosta
$t_n=\sum_{j=0}^{n-1}2^j\chi_{I_j}$
spełnia $0\le t_n\le f$, a jej całka wynosi
\[
\int_{(0,1]}t_n(x)\,\mathrm d\lambda_1(x)
=\sum_{j=0}^{n-1}2^j\lambda_1(I_j)
=\sum_{j=0}^{n-1}2^j2^{-(j+1)}=n/2.
\]
Z definicji całki nieujemnej
$\int_{(0,1]}f(x)\,\mathrm d\lambda_1(x)\ge n/2$
dla każdego $n$, a więc całka jest równa $+\infty$.
Nie trzeba do tego rachunku twierdzenia Beppo Leviego.
\end{przyklad}

\begin{uwaga}[Istnienie całki ze znakiem nie jest tym samym co całkowalność]
Funkcja nieujemna $f(x)=1/x$ na $(0,1]$ ma
określoną całkę równą $+\infty$, ale nie jest
całkowalna w sensie $L^1$. Z kolei dla $f(x)=x$
na $\mathbb R$ obie całki części dodatniej i ujemnej
są nieskończone, a więc całka ze znakiem nie jest
nawet określona. Nie należy utożsamiać trzech sytuacji:
całki skończonej, całki nieskończonej i wyrażenia
nieokreślonego.
\end{uwaga}

\section{Wzajemny stosunek całki i mierzalności}

\begin{przyklad}[Funkcja Dirichleta]
Na przedziale $[0,1]$ funkcja
$f=\chi_{\mathbb Q\cap[0,1]}$ jest borelowsko
mierzalna i nieujemna. Ponieważ zbiór liczb
wymiernych jest przeliczalny, ma miarę Lebesgue'a
zero, a więc $\int_{[0,1]} f\,\mathrm d\lambda_1(x)=0$.
Jednocześnie funkcja jest nieciągła w każdym punkcie,
więc nie jest całkowalna w sensie Riemanna.
\end{przyklad}

\begin{przyklad}[Całkowalność na przestrzeni o nieskończonej mierze]
Na $\mathbb R$ z miarą Lebesgue'a funkcja
$f(x)=1/(1+x^2)$ jest nieujemna i całkowalna, choć
$\lambda_1(\mathbb R)=\infty$. Istotnie, na $[-1,1]$
mamy $f\le1$. Dla $n\ge1$ i $x\in[n,n+1)$ mamy
$f(x)\le n^{-2}$; analogicznie na ujemnej półosi.
Dla dowolnej funkcji prostej $0\le s\le f$,
przeliczalna addytywność miary na poziomicach $s$
pozwala rozłożyć jej całkę na rozłączne przedziały
$[-1,1]$, $[n,n+1)$ i $(-(n+1),-n]$, $n\ge1$.
Na każdym przedziale całka $s$ jest nie większa
od iloczynu jego długości i odpowiednio $1$ albo
$n^{-2}$, więc
$\int_{\mathbb R}s\,\mathrm d\lambda_1(x)
\le2+2\sum_{n\ge1}n^{-2}<\infty$.
Ta nierówność jest niezależna od wyboru $s$;
biorąc supremum w definicji całki, otrzymujemy
$\int_{\mathbb R} f(x)\,\mathrm d\lambda_1(x)<\infty$. Dokładną wartość $\pi$ wyprowadzimy
po omówieniu twierdzenia o zamianie zmiennych.
\end{przyklad}

\begin{uwaga}[Jeszcze nie wszystko zostało dowiedzione]
Na tym etapie znamy monotoniczność, jednorodność,
całki funkcji prostych i kryteria całkowalności.
Nie używaliśmy bez dowodu równości
$\int_X(f+g)=\int_X f+\int_X g$ dla \emph{dowolnych}
nieujemnych funkcji mierzalnych. Tę fundamentalną
własność oraz równość $\int_X\lim f_n=\lim\int_X f_n$
w odpowiednich warunkach wyprowadzimy wprost
z definicji w kolejnym rozdziale.
\end{uwaga}
